\section{Method}
\label{sec:method}

\paragraph{Residual updates as a depth-indexed sequence.} Consider a pre-norm
Transformer with $L$ blocks acting on a hidden state $\hh_\ell \in \mathbb{R}^{T
\times d}$ for $T$ tokens of width $d$. Each block contributes an additive
residual update
\begin{equation}
\hh_{\ell+1} = \hh_\ell + \dd_\ell,
\qquad
\dd_\ell = F_\ell(\hh_\ell),
\label{eq:residual}
\end{equation}
where $F_\ell$ is the attention-plus-MLP sublayer stack. Skipping block $\ell$
means producing an estimate $\widehat{\dd}_\ell$ \emph{without} evaluating
$F_\ell$, then propagating $\widehat{\hh}_{\ell+1} = \hh_\ell +
\widehat{\dd}_\ell$. Seen this way, every layer-skipping method is a
\emph{predictor} of a missing update, and methods differ only in what they
predict.

\paragraph{Plain skipping is a zero-order predictor.} The conventional skip
operator sets
\begin{equation}
\widehat{\dd}^{\,\mathrm{skip}}_\ell = 0 .
\label{eq:plainskip}
\end{equation}
It is the unique estimator that uses no information at all. That is what makes it
cheap, and it is also what makes it wasteful: by the time the model reaches layer
$\ell$ it has already computed $\dd_{\ell-1}, \dd_{\ell-2}, \dots$, all of which
sit in the residual stream and all of which \cref{eq:plainskip} discards. A
second, equally free estimator is available --- reuse the previous update,
$\widehat{\dd}_\ell = \dd_{\ell-1}$ --- and we carry this \textbf{Copy Update}
baseline throughout, because any gain \mname shows over it is a gain from
\emph{fitting}, not merely from being nonzero.

\paragraph{\mname: autoregression across depth.} We model the update sequence as
locally autoregressive in depth. \mname(1) forecasts the missing update from its
immediate predecessor with a single fitted scalar $\alpha_\ell$,
\begin{equation}
\widehat{\dd}_\ell = \alpha_\ell \dd_{\ell-1},
\qquad
\widehat{\hh}_{\ell+1} = \hh_\ell + \alpha_\ell \dd_{\ell-1},
\label{eq:ar1}
\end{equation}
and \mname(2) widens the window to two taps,
\begin{equation}
\widehat{\dd}_\ell = \alpha_\ell \dd_{\ell-1} + \beta_\ell \dd_{\ell-2}.
\label{eq:ar2}
\end{equation}
The inductive bias is momentum: the model has been refining its representation in
a consistent direction, and the next refinement most likely continues it. Note
that $\alpha_\ell = 0$ recovers plain skipping exactly and $\alpha_\ell = 1$
recovers Copy Update, so \mname strictly generalizes both and calibration can
fall back to either when the data demand it.

\paragraph{Closed-form calibration, no gradients.} Coefficients are fitted by
least squares against the true update on a small unlabeled calibration set. For
\mname(1) the solution is a ratio of Frobenius inner products,
\begin{equation}
\alpha^{\star}_\ell =
\frac{\langle \dd_{\ell-1}, \dd_\ell \rangle_F}
     {\|\dd_{\ell-1}\|_F^2 + \epsilon},
\label{eq:fit}
\end{equation}
and for \mname(2) the pair $(\alpha_\ell, \beta_\ell)$ solves a $2 \times 2$
ridge system built from the Gram matrix of $\{\dd_{\ell-1}, \dd_{\ell-2}\}$. Both
are obtained by streaming a few calibration sequences through the dense model
once and accumulating scalar statistics in FP64: no hidden states are stored, no
gradients are taken, and recalibrating an entire model costs a single forward
pass over 16--\ph{32} sequences. Padding tokens are excluded from the
accumulators.

\paragraph{Predictability is measurable before skipping.} The same statistics
yield a held-out score that says, per layer, how much of the update a predictor
explains:
\begin{equation}
P_\ell = 1 -
\frac{\sum \|\dd_\ell - \widehat{\dd}_\ell\|_2^2}
     {\sum \|\dd_\ell\|_2^2}.
\label{eq:pscore}
\end{equation}
$P_\ell > 0$ means the predictor explains more update energy than plain skipping,
and $P_\ell$ is exactly the fraction of that energy it recovers. Because $P_\ell$
is computed from the dense model on unlabeled text, it is available \emph{before}
any skipping decision is taken; we use it both to choose which layers to skip and
to anticipate where skipping will succeed. Unless stated otherwise we skip the
$k$ highest-$P_\ell$ layers subject to a non-adjacency constraint, so that no
skipped block consumes a predicted rather than a real update.

\paragraph{Cost.} \mname adds one or two scaled vector additions per skipped
block: $O(Td)$ work against the $O(Td^2)$ the block itself would have cost, and
$O(1)$ stored parameters per layer rather than the $O(d^2)$ of a fitted linear
map. The predicted update is assembled from tensors already resident in the
residual stream, so the only extra memory traffic is retaining $\dd_{\ell-1}$
(and $\dd_{\ell-2}$ for \mname(2)) for one further layer. The block-compute
reduction of skipping therefore survives essentially intact.
